\section{Methods}
\subsection{Data sources and cohort construction}
The analysis used the 48 half-hour, two-lead Holter recordings provided by the MIT-BIH Arrhythmia Database and harmonised every beat label under the Association for the Advancement of Medical Instrumentation (AAMI) EC57 taxonomy to enable comparison with recent open benchmarks.\cite{Goldberger2000PhysioBank,luz2016ecg} In this mapping the five diagnostic classes are: N (normal and bundle branch beats: N, L, R, e, j, n, B), S (supraventricular ectopic beats: A, a, J, S), V (ventricular ectopic beats: V, E), F (fusion beats: F), and Q (unclassifiable and noisy beats: /, f, Q, ?). Following common AAMI-oriented practice, we retained the subset of 45 records without paced beats for model development by excluding the paced-beat records 102, 104, and 107 from the original corpus.\cite{moody2001impact} The remaining record identifiers 100--124 and 200--234 therefore constitute the AAMI-compliant cohort. From these 45 records, seven subjects (106, 207, 208, 220, 228, 231, 233; 15.6\% of records, 15{,}573 beats) were selected as a strictly held-out test cohort using a record-level random split (\texttt{test\_size}=0.15, random seed~42), while the remaining 38 records were allocated to the training/validation pool. Five-fold cross-validation was performed across these 38 development records before retraining the final model on their union and evaluating it once on the held-out test subjects. This scheme preserves inter-patient independence between development and evaluation sets and ensures that both 1xx and 2xx cohorts---with differing arrhythmia burdens---are represented in the test split, producing more conservative but clinically realistic performance estimates than beat-wise data splitting.

\subsection{Signal conditioning and representation learning}
Signals were left in the form distributed through PhysioNet to keep the benchmark comparable, so no additional band-pass or baseline wander filtering was applied beyond the acquisition settings. R-peak locations and beat annotations were parsed from the PhysioNet \texttt{.atr} files, after which 187-sample windows centered on each fiducial were extracted from the first ECG channel and concatenated into three-beat sequences, i.e., every triple-beat segment is constructed exclusively from lead~1, matching the public preprocessing script (\texttt{preprocess\_data.py}). Each window was transformed into a complex Morlet wavelet scalogram spanning 32 integer scales (1--32), and the magnitude component was retained to form 32\,$\times$\,187 time--frequency images that emphasize multiresolution morphology while preserving temporal alignment.\cite{addison2005wavelet,torrence1998practical,mallat1989wavelet} The scalograms plus metadata were written to HDF5 containers (\texttt{preprocessed\_data\_h5\_raw/}). Amplitude normalization was deferred to the dataset-construction stage so that fold-specific statistics could be applied consistently by \texttt{DataLoader.py}. No augmentation such as jittering, polarity inversions, or synthetic resampling was introduced in the experiments reported here. We restrict analysis to the first ECG lead to simplify deployment and align with prior MIT-BIH benchmarks that operate on single-lead or lead-aggregated representations.\cite{kiranyaz2015personalized,luz2016ecg}

\subsection{Neural architecture and training procedure}
The classifier combines a shallow convolutional stem with Conformer blocks for contextual modelling.\cite{gulati2020conformer} Hyperband explored embedding sizes (64 or 128), head counts (4 or 8), kernel widths (7 or 11), dropout rates between 0.1 and 0.3, and one or two Conformer blocks, and selected a 128-dimensional, 8-head, 11-sample configuration for the final model. Every model was trained with AdamW (learning rate $1\times10^{-5}$) for 20 epochs using class-weighted sparse categorical cross-entropy on NVIDIA-class GPUs. Comparator models include an attention-only variant that removes the convolution branch and a CNN-LSTM baseline that substitutes recurrent units for the Conformer encoder, enabling controlled ablations of architectural components.

\subsection{Evaluation metrics and statistical analysis}
Model selection prioritised macro-averaged F1 to balance sensitivity across arrhythmia types, while accuracy, weighted F1, and one-vs-rest area under the receiver operating characteristic curve (AUROC) provided complementary summaries. \texttt{Evaluator.py} produced the held-out metrics, confusion matrix, and ROC as well as precision--recall plots directly from the stored predictions, ensuring that every number in the manuscript can be regenerated through a single command. Five-fold cross-validation on the 38 training records yielded mean and standard-deviation accuracy values via the automated \texttt{run\_kfold\_evaluation.py} routine, offering a coarse measure of subject-level variability. For calibration, we additionally compute per-class expected calibration error (ECE) and Brier scores from the stored test-set probabilities using 15 equally spaced confidence bins between 0 and~1; ECE is defined as the weighted average absolute difference between bin-wise accuracy and mean confidence, and the Brier score as the mean squared difference between predicted probabilities and one-hot labels.
Ninety-five percent confidence intervals for test-set accuracy, macro-F1, weighted F1, and AUCs in \autoref{tab:test-metrics} were obtained via nonparametric bootstrap resampling of the 15{,}573 test beats with replacement; this procedure operates at the beat level and does not explicitly model within-patient correlation.

\subsection{Software and hardware environment}
All experiments were executed in Python~3.12.10 with TensorFlow~2.20.0 on Ubuntu running inside the Windows Subsystem for Linux. Core dependencies included \texttt{wfdb}~4.1.2 for waveform access, \texttt{PyWavelets}~1.6.0 for continuous wavelet transforms, \texttt{seaborn}~0.13.2 for plotting, and standard scientific Python libraries (NumPy, SciPy, pandas, and scikit-learn) at contemporary 2025 releases, matching the versions documented in the public repository. Training used an NVIDIA GeForce GTX~1660~Ti GPU (6~GB VRAM), an Intel Core~i7 CPU, and 16~GB RAM, with CUDA~12.2 and cuDNN~9.0 for hardware acceleration.
Reproducibility is enforced through fixed random seeds, version-pinned \texttt{requirements.txt}, and archived configuration files for every experiment, enabling reviewers to validate the reported metrics using the provided scripts.
